\documentclass[12pt, letterpaper]{article}

\title{CS302: Error Handling \color{red}\textbf{[KEY]}\normalcolor}
\author{Amiel L. Iliesi}
\date{2026-03-30}

% packages
\usepackage[hidelinks]{hyperref}
\usepackage[x11names]{xcolor}
\usepackage{minted}

% new commands
\newcommand{\question}[1]{\textbf{#1}\par}
\newcommand{\answer}[1]{\color{Green4}\textbf{#1}\normalcolor}

% new environments
\newenvironment{answerblock}{\color{Green4}}{\normalcolor\vspace{1em}}

% global settings
\setlength{\parindent}{0pt}
\setlength{\parskip}{0.5em}
\usemintedstyle{monokai}
\setminted{bgcolor=darkgray}

\begin{document}

\maketitle

\tableofcontents

\pagebreak
\section{C Style: Error Codes}

\question{How do C-style programs and functions typically indicate errors?}

\begin{answerblock}
Via return type, typically through error codes or special return values.
\end{answerblock}

\question{What are some pros and cons to this approach? Why and how does modern
error handling differ?}

\begin{answerblock}
Some pros are that it requires very little runtime overhead, since it's just
using the function return infrastructure. It also results in predictable
program flow, as there isn't remote error prompted stack unwinding.
Finally--and arguably a con--you aren't forced to resolve the error, if it is
trivial or non-fatal.

Some cons are that there is no standard for this return method, so you rely on
both: that the programmer that returned the code--returned the code correctly,
and that the documented codes are returned in every case. Blind trust is a poor
error handling paradigm. Also what was listed as a potential pro, that as a
recipient of the code--your handling of the error is not enforced--also produces
poor adherence to error handling. This will lead to untested, unhandled errors
that are silently ignored.
\end{answerblock}

\pagebreak
\section{C++}

\subsection{\texttt{stdexcept} and \texttt{throw}}

\begin{answerblock}
\texttt{stdexcept} is the C++ library that has a few standardized error classes
for us to use for our return types. Throw is what we call within a function to
start the error return process. It tells the function to stop and unwind the
stack from that given line, and carry that error value with it.
\end{answerblock}

\question{What is \texttt{stdexcept}? Give a few examples.}
\begin{answerblock}
\texttt{stdexcept} comes with some standardized error classes that we can
inherit from, or use directly. Because they are standardized, their intent is
well understood. For example, an error class that \texttt{stdexcept} has is
\texttt{std::logic\_error} which indicates that a runtime logic error has
occurred. Another class is \texttt{std::out\_of\_range}, which indicates that
there was an attempted access out of a given range.
\end{answerblock}

\question{What does the \texttt{throw} keyword do? How is this different from
\texttt{return}?}

\begin{answerblock}
\texttt{return} pops from the stack \emph{once}, and returns a definite,
singular type. \texttt{throw}--on the other hand--keeps popping until it reaches
a \texttt{try}/\texttt{catch} block, and enters the matching except block,
\emph{or} until it exits the \texttt{main} call and aborts with an error status.

In short, the \texttt{return} and \texttt{throw} follow separate paths, with
slightly different behaviors.
\end{answerblock}

\begin{minipage}{\textwidth}
\question{What if \texttt{stdexcept} doesn't have what we need? How would you
write something more specific to your needs?}
\begin{answerblock}
There are three good options.

\begin{enumerate}
    \item Inherit a new class from one of the base classes defined in the
    \texttt{stdexcept} library.
    \item Inherit a new class from the base exception class defined in the
    \texttt{exception} library.
    \item If both of the above don't suit your needs, create a custom
    \texttt{class} or \texttt{struct} to throw.
\end{enumerate}
\end{answerblock}
\end{minipage}

\subsection{\texttt{try}/\texttt{catch}}

\question{What are \texttt{try} and \texttt{catch}?}

\begin{answerblock}
If you can expect a piece of code to fail and throw an error, it should be
wrapped in a \texttt{try}/\texttt{catch} block. The \texttt{try}/\texttt{catch}
ensures predictable flow in the event of an error, IE "I have seen that the
error occured, and now to keep running, do this instead".

For example, if you are doing a division with floating point numbers, you
should wrap the computing block or function in a \texttt{try}/\texttt{catch},
to handle a divide-by-zero.
\end{answerblock}

\question{What are the \emph{three} parameter formats you can specify for
\texttt{catch}?}

\begin{answerblock}
\begin{itemize}
    \item catch based on type, anonymously:
    \mintinline{c++}|catch(TYPE&) {}|
    \item catch based on type:
    \mintinline{c++}|catch(TYPE& err_obj_name) {}|
    \item catch all:
    \mintinline{c++}|catch(...) {}|
\end{itemize}
\end{answerblock}

\subsection{Errors as Return Values, and \texttt{std::expected}}

\textbf{NOTE: you may not need to know this for class, but it's good to know
for modern error handling practices.}

\question{What is an \emph{error as return value}? Describe the pros and cons
to this approach. Also show a simple example of {\lq{throwing}\rq}  an error value with
\texttt{std::expected}.}

\begin{answerblock}
For exceptions, throwing is much easier and more readable
\emph{at point of error}. It makes the throw site more readable, but it comes
with a host of negatives. Because it can throw \emph{any type}, the domain of
values the reader has to deal with at the \texttt{catch} block are vague and
frustrating. This leads to a messy, nested, obfuscated execute block at the
\texttt{try}/\texttt{catch}, lovingly referred to as
\texttt{try}/\texttt{catch} \emph{hell}. As a nice bonus, there's a runtime
cost for using exceptions.

Using errors as values on the other hand--\emph{if used correctly}--fixes a lot
of these issues. It's why \textbf{safe} languages like Rust adopted this style
for their exception handling. You form a compound type that can either be the
value you expect (like an int) \emph{or}, it can be an unexpected error value.
This means that the user is forced to deal with the error, \emph{and} that the
error type is known to the caller beforehand as well. Finally, since it uses
the \texttt{return} infrastructure, it is also fast.

In short, the convenience and flexibility of exceptions mean that they,
ironically, become hard to structure and decipher from the receiver's
perspective. Having rigid errors as values makes things predictable and
readable.

And here's an error as value example with \texttt{std::expected}:

\begin{minipage}{\textwidth}
\begin{minted}{cpp}
enum class Error {
    divide_by_zero,
    NaN
};

auto f(float a, float b) -> std::expected<float, Error>
{
    if std::isnan(a) or std::isnan(b) {
        return std::unexpected(Error::NaN);
    }
    else if (b == 0.0) {
        return std::unexpected(Error::divide_by_zero);
    }

    return a/b;
}

int main() {
    float a = 5.0;
    float b = 3.0;

    auto res = f(a, b);

    if (res.has_value()) {
        std::println("{}/{}={}", a, b, *res);
    }
    else if (res.error() == Error::divide_by_zero) {
        std::println(stderr, "Cannot divide by zero");
    }
    else if (res.error() == Error::NaN) {
        std::println(stderr, "Cannot do division with NaN");
    }
}
\end{minted}
\end{minipage}
\end{answerblock}

\pagebreak
\section{Python}

\subsection{Raise}

\question{What are the differences between Python's \texttt{raise}, and C++'s
\newline\texttt{throw}?}

\begin{answerblock}
Besides the differences in syntax from Python to C++, they are equivalent.
\end{answerblock}

\question{What types of exceptions can you throw in Python? How can you define
more--if needed?}

\begin{answerblock}
Same as in C++, inherit from a base or derived exception type. EG:

\begin{minted}{python}
class customException(Exception):
    pass

raise customException
\end{minted}
\end{answerblock}

\subsection{\texttt{try}/\texttt{except}/\texttt{else}/\texttt{finally}}

\question{Explain these four keywords.}

\begin{answerblock}
\begin{itemize}
    \item try: code to execute.
    \item except: exception to handle, same as C++'s catch.
    \item else: code to execute if \textbf{no} exceptions were raised.
    \item finally: code to run afterwards, unconditionally. It runs before the
    block is popped from the stack, IE before any returns or breaks that may
    be present.
\end{itemize}
\end{answerblock}

\end{document}
